\subsection{Computational Methods for Fixed Daily Sample Size}

A pitfall of treating the data in its original form, as opposed to calculating aggregates at each time step, is the relative strain this can put on computation and memory. We can leverage certain aspects of the model to make choices in our computational methods that greatly improve computing time and reduce memory requirements.

\subsection{Inverting $\Omega$}
The biggest computational bottleneck in this procedure is the calculation and inversion of $\Omega$, a necessary step in estimating parameters and calculating the log-likelihood. We can expedite these processes with expressions for $\Omega$ that relate its spectral decomposition to that of $\Gamma$.

We begin by recalling the expression for $\Omega$ and defining $P$ and $\Lambda$ as the components of the spectral decomposition of $\mathbf{V} = \sigma_\eta^2 \GG$:

$$\Omega = X_t\Gamma X_t' + \sigma_\varepsilon^2I_N, \qquad  \Gamma = \sigma_\eta^2P \Lambda P'$$

Because $\Gamma$ is positive definite, all eigenvalues $\{\lambda_i\}_{i=1}^T$ are real-valued and strictly positive. Let $\{\mathbf{p}_i\}_{i=1}^T$ be the corresponding set of eigenvectors. If daily sample sizes are a constant $n$, then $X_t\mathbf{p}_i$ is an eigenvector of $X_t\Gamma X_t'$ with corresponding eigenvalue $n\lambda_i$:

$$ (X_t \Gamma X_t')(X_t\mathbf{p}_i) = (X_t\Gamma)(X_t'X_t)\mathbf{p}_i = X_t \Gamma n {I}_T\mathbf{p}_i = nX_t\Gamma \mathbf{p}_i = n\lambda (X_t \mathbf{p}_i) $$

Given that $\mathbf{p}_i$ is a normal vector and $X_t\mathbf{p}_i$ has $n$ copies of each element of $\mathbf{p}_i$, then it necessarily follows that $||X_t\mathbf{p}_i||^2 = n$. We can then obtain a normal eigenvector with $n^{-1/2}X_t\mathbf{p}_i$ and verify that the resulting set $\{n^{-1/2}X_t\mathbf{p}_i\}$ is orthonormal:

$$ \langle n^{-1/2}X_t\mathbf{p}_i, n^{-1/2}X_tv_j \rangle = n^{-1/2}X_t\mathbf{p}_i v_j'X_t'n^{-1/2} = n^{-1/2}X_t\mathbf{0}X_t'n^{-1/2} = \mathbf{0}, \qquad i\neq j$$

Now, $X_t\Gamma X_t'$ is an $N \times N$ matrix expressable as a $T \times T$ block matrix, for which the $i,j$ block is $\gamma_\uu(i-j) J_n$, with $J_n$ defined as an $n \times n$ matrix of ones. As a result, $N-T$ rows of $X_t\Gamma X_t'$ are simply copies of the other $T$ rows, meaning $\{n^{1/2}\lambda_i\}_{i=1}^T$ are the only nonzero eigenvalues of $X_t\Gamma X_t'$. Because $X_t\Gamma X_t'$ is symmetric, any vector in its null space is orthogonal to all eigenvectors corresponding to nonzero eigenvalues. 

This means we just need a $N \times (N-T)$ matrix whose columns form an orthonormal basis for the null space of $X_t\Gamma X_t'$. Before defining said matrix, let's call it $Q$ and motivate its derivation by detailing its use in cutting down computing time:

$$ X_t \Gamma X_t' = \left[\begin{array}{cc}
   n^{-1/2} X_t P & Q  
\end{array}\right] \left[\begin{array}{cc}
    n^{1/2}\sigma_\eta^2\Lambda & \mathbf{0} \\
    \mathbf{0} & \mathbf{0}
\end{array}\right] \left[\begin{array}{cc}
    n^{-1/2}P'X_t'  \\
    Q' 
\end{array}\right] $$

Because the columns of $n^{-1/2}X_tP$ and $Q$ are all orthonormal and $\Gamma$ is expressable as $\sigma_\eta^2 V$, we can now give the following spectral decomposition for $\Omega$:

% $$ \Omega = X_t\Gamma X_t' + \sigma_\varepsilon^2I_N = \left[\begin{array}{cc}
%    n^{-1/2} X_t P & Q  
% \end{array}\right] \left[\begin{array}{cc}
%     n^{1/2}\sigma_\eta^2\Lambda & \mathbf{0} \\
%     \mathbf{0} & \mathbf{0}
% \end{array}\right] \left[\begin{array}{cc}
%     n^{-1/2}P'X_t'  \\
%     Q' 
% \end{array}\right]  + \left[\begin{array}{cc}
%    n^{-1/2} X_t P & Q  
% \end{array}\right] \left[\sigma_\varepsilon^2I\right] \left[\begin{array}{cc}
%     n^{-1/2}P'X_t'  \\
%     Q' 
% \end{array}\right] $$

\begin{equation}
\label{eq:eigen_omega}
    \Omega = \left[\begin{array}{cc}
   n^{-1/2} X_t P & Q 
\end{array}\right] \left[\begin{array}{cc}
    n^{1/2}\sigma_\eta^2\Lambda + \sigma_\varepsilon^2I_T & \mathbf{0} \\
    \mathbf{0} & \sigma_\varepsilon^2I_N
\end{array}\right] \left[\begin{array}{cc}
    n^{-1/2}P'X_t'  \\
    Q' 
\end{array}\right]  
\end{equation} 

We now derive an expression for the inverse of $\Omega$, for which we only have to invert a $T \times T$ diagonal matrix:

$$ \Omega^{-1} =  \left[\begin{array}{cc}
   n^{-1/2} X_t P & Q
\end{array}\right] \left[\begin{array}{cc}
    (n^{1/2}\sigma_\eta^2\Lambda + \sigma_\varepsilon^2I_T)^{-1} & \mathbf{0} \\
    \mathbf{0} & \sigma_\varepsilon^{-2} I_N
\end{array}\right] \left[\begin{array}{cc}
    n^{-1/2}P'X_t'  \\
    Q' 
\end{array}\right]  $$

$$\Omega^{-1} =  
   n^{-1} X_t P (n^{1/2}\sigma_\eta^2\Lambda + \sigma_\varepsilon^2I_T)^{-1}P'X_t' + \sigma_\varepsilon^{-2}Q
   Q'$$

Having motivated its derivation, we now define $Q$. Let $v$ be a vector in the null space of $X_t\Gamma X_t'$. By definition $v$ must be orthogonal to every row vector in $X_t\Gamma X_t'$, and due to symmetry every column vector. If orthogonal to every column vector, $v$ must be orthogonal to any vectors that form a basis for the column space, namely the eigenvectors forming the columns of $n^{-1/2}X_tP$. If we're able to find $N-T$ vectors that are orthonormal to each other and the eigenvectors previously defined, we will necessarily have constructed a basis for the null space of $X_t\Gamma X_t'$, and in doing so obtained the missing $N-T$ eigenvectors thereof. Let us first note that the matrix $X_tP$ can be expressed as the following block matrix:


$$X_tP =\left[\begin{array}{c}
    \vdots  \\
    P_t \\
    \vdots 
\end{array}\right], \qquad P_t \in \mathbb{R}^{T\times T}\qquad (P_t)_{i,j} = (\mathbf{p}_j)_t$$

Noting that the entries of $P_t$ are independent of $i$, the column space of $P_t$ is equal to that of $\mathbf{1}_n$. If we define $\{\mathbf{q}_i\}_{i=1}^{n-1}$ to be an orthonormal basis for $\mathcal{N}(\mathbf{1}_n)$ and $\mathbf{Q}_t$ to be the $n \times (n-1)$ matrix with column vectors $\{\mathbf{q}_i\}$, we then define the matrix $\mathbf{Q}$ as follows:

$$ \mathbf{Q} = \left[\begin{array}{ccc}
    \ddots &  &  \mathbf{0}\\
     & \mathbf{Q}_j & \\
    \mathbf{0} &  & \ddots
\end{array}\right], \qquad j=1,...,T   $$





The columns of $\mathbf{Q}_j$ are necessarily orthogonal to those of $\mathbf{Q}_i$ for $i\neq j$ since no pair of corresponding entries are both nonzero, and the columns of $\mathbf{Q}_j$ are orthogonal to each other since the columns of $\mathbf{Q}_{j,t}$ are orthogonal and all other entries are zero. It is for this same reason that the columns of $\mathbf{Q}_i$ are of magnitude $1$. We can now express $\mathbf{QQ'}$ as the computationally cheap block diagonal matrix:

$$ QQ' = \left[ \begin{array}{ccc}
    \ddots &  & \mathbf{0} \\
      & \mathbf{Q}_j\mathbf{Q}_j' &  \\
     \mathbf{0} & & \ddots
\end{array} \right] $$

Equipped with this expression for $\Omega^{-1}$, we demonstrate how it expedites the various parts of the estimation algorithm.

\subsection{Estimating $\uu$}
Though it is no longer used for parameter estimation, it is still valuable to estimate $\uu$ for the purposes of identifying any short term trends that would otherwise be lost if only fitting a linear trend. Recall $\hat{\uu}$ is given by $\GG X_t' \Omega^{-1}(\yy - X_f\hat{\beta})$, which is now expressible as:

$$ \hat{\uu} = \sigma_\eta^2(P\Lambda P'X_t')(n^{-1} X_t P (n\sigma_\eta^2\Lambda + \sigma_\varepsilon^2I_T)^{-1}P'X_t' + \sigma_\varepsilon^{-2}Q
   Q')(\yy-X_f\hat{\beta}) $$
   
$$ \hat{\uu} = n^{-1}\sigma_\eta^2P\Lambda (n\sigma_\eta^2\Lambda + \sigma_\varepsilon^2I_T)^{-1}P'X_t'(\yy-X_f\beta) + \sigma_\varepsilon^{-2}\sigma_\eta^2(P\Lambda P'X_t')Q
   Q')(\yy - X_f\hat{\beta}) $$

By virtue of the orthogonality of the column vectors of $X_tP$ and those of $Q$, the above expression simplifies to:

$$ \hat{\uu} = \sigma_\eta^2P\Lambda (n\sigma_\eta^2\Lambda + \sigma_\varepsilon^2I_T)^{-1}P'(\bar{\yy} - n^{-1}X_t'X_f\beta) $$

The punchline of this expression is that what was originally an expression involving the calculation of the inverse of a $nT \times nT$ matrix and multiplication thereof with matrices of dimension $nT \times T$ has now been re-expressed as one involving matrices of no greater dimension than $T \times T$. With this expression we can reduce the calculation of $\hat{\uu}$ from $O(N^3)$ to $O(T^3)$, which differs by a factor of $n^3$.

\subsection{Calculating $\hat{\beta}_\text{OLS}$}

Just as the calculation and storage of $\Omega$ and its inverse for the purposes of estimating $\uu$ can strain memory, so too is it a problem in the calculation of $\hat{\beta}_{GLS}$, defined as: 

    $$\hat{\beta}_{GLS} = (X_f'\Omega^{-1}X_f)^{-1} X_f'\Omega^{-1}\yy$$
We can change the order of operations to minimize computing time:
\begin{equation}
    X_f'\Omega^{-1}X_f = X_f' (n^{-1} X_t P (n^{1/2}\sigma_\eta^2\Lambda + \sigma_\varepsilon^2I_T)^{-1}P'X_t' + \sigma_\varepsilon^{-2}Q
   Q') X_f
\end{equation}
We now distribute $X_f'$ and $X_f$ on the right hand side.

\begin{equation}
    X_f'\Omega^{-1}X_f =  n^{-1}X_f' X_t P (n^{1/2}\sigma_\eta^2\Lambda + \sigma_\varepsilon^2I_T)^{-1}P'X_t'X_f + \sigma_\varepsilon^{-2} X_f'Q
   Q' X_f
\end{equation}

If we then prioritize calculating $X_f'X_t$ and $X_f'Q$, which are respectively matrices of dimension $k \times T$ and $k \times (N-T)$, we then greatly reduce the maximum size of matrix in the expression. Inverting $X_f'\Omega^{-1}X_f$. 

